% Options for packages loaded elsewhere
\PassOptionsToPackage{unicode}{hyperref}
\PassOptionsToPackage{hyphens}{url}
\documentclass[
]{article}
\usepackage{xcolor}
\usepackage{amsmath,amssymb}
\setcounter{secnumdepth}{-\maxdimen} % remove section numbering
\usepackage{iftex}
\ifPDFTeX
  \usepackage[T1]{fontenc}
  \usepackage[utf8]{inputenc}
  \usepackage{textcomp} % provide euro and other symbols
\else % if luatex or xetex
  \usepackage{unicode-math} % this also loads fontspec
  \defaultfontfeatures{Scale=MatchLowercase}
  \defaultfontfeatures[\rmfamily]{Ligatures=TeX,Scale=1}
\fi
\usepackage{lmodern}
\ifPDFTeX\else
  % xetex/luatex font selection
\fi
% Use upquote if available, for straight quotes in verbatim environments
\IfFileExists{upquote.sty}{\usepackage{upquote}}{}
\IfFileExists{microtype.sty}{% use microtype if available
  \usepackage[]{microtype}
  \UseMicrotypeSet[protrusion]{basicmath} % disable protrusion for tt fonts
}{}
\makeatletter
\@ifundefined{KOMAClassName}{% if non-KOMA class
  \IfFileExists{parskip.sty}{%
    \usepackage{parskip}
  }{% else
    \setlength{\parindent}{0pt}
    \setlength{\parskip}{6pt plus 2pt minus 1pt}}
}{% if KOMA class
  \KOMAoptions{parskip=half}}
\makeatother
\setlength{\emergencystretch}{3em} % prevent overfull lines
\providecommand{\tightlist}{%
  \setlength{\itemsep}{0pt}\setlength{\parskip}{0pt}}
\usepackage{bookmark}
\IfFileExists{xurl.sty}{\usepackage{xurl}}{} % add URL line breaks if available
\urlstyle{same}
\hypersetup{
  pdftitle={The Trinity as a Necessary Structure of a Monistic Modal Ontology},
  hidelinks,
  pdfcreator={LaTeX via pandoc}}

\title{The Trinity as a Necessary Structure of a Monistic Modal Ontology \\ \large A Strictly Formalized Modal-Logical Investigation}
\author{Paul Koop}
\date{}

\begin{document}
\maketitle

\begin{center}
\emph{Past and future are horizons of knowledge. \\ The present is the place of reality.}
\end{center}

\newpage
{
\setcounter{tocdepth}{3}
\tableofcontents
}
\newpage

\section{Preface}\label{preface}

This text is an attempt to demonstrate an ancient metaphysical intuition -- the threefold unity of origin, totality, and self-knowledge -- not as a matter of faith, but as a logical necessity of a consistent ontology.

The investigation is conducted within the modal-logical system S5. It strictly distinguishes between three types of proof: the ontological proof of existence, the consistency proof, and the transcendental argument of necessity. The central thesis is:

> Every possible world that obeys the principles of monism, non-emergence from nothing, modal realism, S5 modality, the real possibility of consciousness, and the reflexivity of complete knowledge must necessarily exhibit the structure of the Trinity.

This treatise does not provide a proof of God in the classical sense. It does not demonstrate a personal or substantial being. It merely shows that, under a defined system of axioms, the threefold structure of Origin, Totality, and Self-Knowledge is unavoidable. It is not an additional entity, but a structural condition.

The text is formal in its approach but written in the language of philosophy -- for truth requires both: the precision of the formula and the breadth of the concept.

\section{Introduction: The Problem of the Three Perspectives}\label{introduction-the-problem-of-the-three-perspectives}

The philosophical tradition knows numerous triads: Being--Thought--Spirit, Ground--Existence--Consummation, Father--Son--Holy Spirit. What they share is the intuition that ultimate reality consists not in a simple unity, nor in a dualistic duality, but in a threefold unity.

The proof presented here takes up this intuition, divests it of its theological garb, and formulates it as an ontological structural condition.

We define three limit-values:

\begin{itemize}
\item \textbf{Origin (U)}: The necessary limit-value from which all possibilities emerge. U is not a first event in time, but the transcendental condition of the possibility of all modalities.
\item \textbf{Totality (T)}: The necessary limit-value of the totality of all realized possibilities. T is not the mere sum of all facts, but the unity of their interconnectedness.
\item \textbf{Self-Knowledge (S)}: The necessary limit-value of an asymptotically complete self-knowledge of the totality, incorporated into reality. S is not an individual consciousness, but the self-referentiality of reality as a whole.
\end{itemize}

It is essential that these three limit-values are not understood as three separate substances, but as three necessary perspectives on the same reality. The proof will show that these perspectives are not only compatible but logically equivalent -- they mutually imply one another.

The Trinity is defined as:

\[
Tr := U \land T \land S
\]

\section{Formal Language and Logical Framework}\label{formal-language-and-logical-framework}

We work in first-order modal logic with identity in the system S5.

\subsection{Modal Logic S5}\label{modal-logic-s5}

\begin{itemize}
\item \textbf{Necessity:} \(\Box p\)
\item \textbf{Possibility:} \(\Diamond p := \neg\Box\neg p\)
\end{itemize}

\textbf{Axioms of S5:}
\begin{itemize}
\item (K) \(\Box(p \rightarrow q) \rightarrow (\Box p \rightarrow \Box q)\)
\item (T) \(\Box p \rightarrow p\)
\item (4) \(\Box p \rightarrow \Box\Box p\)
\item (5) \(\Diamond p \rightarrow \Box\Diamond p\)
\end{itemize}

\textbf{Inference Rules:}
\begin{itemize}
\item \textbf{Modus Ponens:} From \(p\) and \(p \rightarrow q\), infer \(q\).
\item \textbf{Necessitation:} From \(p\) (derivable), infer \(\Box p\).
\end{itemize}

\subsection{Predicate Logic}\label{predicate-logic}

We use classical first-order predicate logic with identity (\(=\)) and the usual quantifiers (\(\forall, \exists\)).

\section{Basic Concepts (Definitions)}\label{basic-concepts-definitions}

All definitions are chosen so as not to already contain the conclusion.

\begin{tabular}{lll}
\textbf{Concept} & \textbf{Definition} & \textbf{Remark} \\
World & \(Welt(w)\) & w is a possible world \\
Possibility & \(Möglich(p) := \Diamond p\) & p is possible \\
Realization & \(Realisiert(w,p)\) & p is realized in world w \\
Totality & \(T := \forall p(\Diamond p \rightarrow \exists w\, Realisiert(w,p))\) & Set of all realized possibilities \\
Consciousness & \(C(x)\) & x is a state of consciousness \\
Self-Consciousness & \(C_s(x) := C(x) \land C(C(x))\) & Consciousness that knows itself \\
Complete Knowledge of T & \(V_T(x)\) & x knows all true statements about T \\
Ground & \(Grund(g,p)\) & g is the ontological ground of p \\
Necessary Existence & \(Nec(x) := \Box \exists x\) & x exists necessarily \\
Fundamental Separation & \(FundamentalGetrennt(x,y)\) & no causal or ontological connection \\
\end{tabular}

\subsection{Origin (U) - Definition}\label{origin-u-definition}

\[
\boxed{U(x) := Grund(x, T) \land \Box \exists x}
\]

\textbf{Explanation:} U is the necessary ground of the totality T. This definition does not presuppose that U exists -- it defines what it means for U to exist. Existence must follow from the axioms.

\section{Axiom System}\label{axiom-system}

All axioms hold necessarily (\(\Box\)).

\subsection{A1: Strict Monism}\label{a1-strict-monism}

There are no fundamentally separate realms of reality.

\[
\boxed{A1 := \Box\neg\exists x\exists y\, FundamentalGetrennt(x,y)}
\]

\textbf{Justification:} If two realms were fundamentally separate, there could be no causal or ontological interaction between them. They could not be part of a unified reality, which would contradict monism. Such a dualism would be incompatible with a complete modal realism.

\subsection{A2: Modal Realism}\label{a2-modal-realism}

Every logically consistent possibility is realized in some possible world.

\[
\boxed{A2 := \forall p(Cons(p) \rightarrow \Diamond Realisiert(p)) \land \forall p(\Diamond p \rightarrow \exists w\, Realisiert(w,p))}
\]

\subsection{A3: No Absolute Nothing}\label{a3-no-absolute-nothing}

An absolute state of nothingness cannot produce any world or possibility.

\[
\boxed{A3 := \Box(N \rightarrow \neg\Diamond Welt)}
\]

Contrapositive:

\[
\Box(\Diamond Welt \rightarrow \neg N)
\]

\subsection{A4: Existence of a Possible Reality}\label{a4-existence-of-a-possible-reality}

There is at least one possible world.

\[
\boxed{A4 := \Diamond\exists w\, Welt(w)}
\]

\subsection{A5: Consciousness as a Real Possibility}\label{a5-consciousness-as-a-real-possibility}

Consciousness is a consistent possibility.

\[
\boxed{A5 := \Diamond C}
\]

\subsection{A6: Origin as an Ontological Embedding Condition}\label{a6-origin-as-an-ontological-embedding-condition}

Every reality that is not necessary in itself requires an ontological ground.

\[
\boxed{A6 := \Box(\neg Nec(p) \rightarrow \exists g\, Grund(g,p))}
\]

\textbf{Explanation:} This is the precise formal version of the principle that contingent entities require a ground. It is weaker than the classical principle of sufficient reason because it applies only to non-necessary entities.

\subsection{A7: Possibility of Complete Knowledge}\label{a7-possibility-of-complete-knowledge}

If consciousness is realized, then complete knowledge of the totality is possible.

\[
\boxed{A7 := \Box(C \rightarrow \Diamond V_T)}
\]

\subsection{A8: Reflexivity of Complete Knowledge}\label{a8-reflexivity-of-complete-knowledge}

A state of complete knowledge knows its own knowledge.

\[
\boxed{A8 := \Box(V_T(x) \rightarrow V_T(V_T(x)))}
\]

\subsection{A9: Identity of Knowledge and Object in Monism}\label{a9-identity-of-knowledge-and-object-in-monism}

If a state of complete knowledge of the totality exists, then the totality must exist as the object of knowledge.

\[
\boxed{A9 := \Box(V_T(x) \rightarrow T)}
\]

\textbf{Explanation:} In monism, there can be no ontological separation between knower and known. Therefore, the existence of \(V_T\) implies the existence of \(T\).

\subsection{A10: Well-Foundedness of the Grounding Relation}\label{a10-well-foundedness-of-the-grounding-relation}

The grounding relation is well-founded; there are no infinite chains of grounds.

\[
\boxed{A10 := \Box\neg\exists infinite\, chain(g_1, g_2, ...)}
\]

\textbf{Explanation:} This axiom prevents an infinite regress and ensures the existence of a final, ungrounded ground -- precisely U.

\section{Proof: Reductio ad absurdum}\label{proof-reductio-ad-absurdum}

\subsection{Assumption}\label{assumption}

We assume that the Trinity does not exist:

\[
\neg Tr \equiv \neg(U \land T \land S)
\]

By de Morgan:

\[
\neg U \lor \neg T \lor \neg S
\]

We must show that each of the three cases leads to a contradiction.

\subsection{Theorem 1: \(\Box(T \rightarrow U)\)}\label{theorem-1-boxt-rightarrow-u}

\textbf{Proof:}

1. From the definition of \(T\):  
   \[
   T \rightarrow \exists p(\Diamond p \land \exists w\, Realisiert(w,p)) \rightarrow \exists w\, Welt(w)
   \]

2. From A3 (contrapositive):  
   \[
   \Box(\Diamond Welt \rightarrow \neg N)
   \]

3. Since a world exists, it is not absolute nothingness. Therefore, it must have a ground:  
   \[
   \exists w\, Welt(w) \rightarrow \exists g\, Grund(g, \exists w\, Welt(w))
   \]

4. If this ground is not itself necessary, then by A6 it is grounded by a further ground. By A10, there is no infinite regression. Therefore, there must be a necessary ground.

5. This necessary ground is precisely \(U\):  
   \[
   U := Grund(g, T) \land \Box \exists g
   \]

6. Therefore:  
   \[
   T \rightarrow U
   \]

7. By Necessitation:  
   \[
   \boxed{\Box(T \rightarrow U)}
   \]

\subsection{Theorem 2: \(\Box(U \rightarrow S)\)}\label{theorem-2-boxu-rightarrow-s}

\textbf{Proof:}

1. From A5: \(\Diamond C\).

2. With A2: \(\Diamond C \rightarrow \exists w\, Realisiert(w,C)\).  
   Thus, a world with consciousness exists.

3. From A7: \(\Box(C \rightarrow \Diamond V_T)\).  
   Thus: \(\exists w\, Realisiert(w,C) \rightarrow \Diamond V_T\).

4. With A2: \(\Diamond V_T \rightarrow \exists w\, Realisiert(w,V_T)\).

5. From A8: \(\Box(V_T(x) \rightarrow V_T(V_T(x)))\).  
   Thus, in this world, \(S\) (self-knowledge) exists.

6. Since \(U\) is the ground of all possibilities (definition of U), we have:  
   \[
   U \rightarrow \forall p(\Diamond p \rightarrow \exists w\, Realisiert(w,p))
   \]

7. Therefore:  
   \[
   U \rightarrow S
   \]

8. By Necessitation:  
   \[
   \boxed{\Box(U \rightarrow S)}
   \]

\textbf{Check:} The implication in step 6 is correct: If U is the ground of T, and T encompasses all realized possibilities, then U is also the ground of each individual realized possibility.

\subsection{Theorem 3: \(\Box(S \rightarrow T)\)}\label{theorem-3-boxs-rightarrow-t}

\textbf{Proof:}

1. From the definition of \(S\):  
   \[
   S \rightarrow \exists x\, V_T(x)
   \]

2. From A9: \(\Box(V_T(x) \rightarrow T)\).  
   Thus: \(\exists x\, V_T(x) \rightarrow T\).

3. Therefore:  
   \[
   S \rightarrow T
   \]

4. By Necessitation:  
   \[
   \boxed{\Box(S \rightarrow T)}
   \]

\subsection{Synthetic Conclusion}\label{synthetic-conclusion}

From the three theorems, it follows:

\[
\Box(T \rightarrow U), \quad \Box(U \rightarrow S), \quad \Box(S \rightarrow T)
\]

In S5, this implies:

\[
\Box(U \leftrightarrow T), \quad \Box(T \leftrightarrow S), \quad \Box(S \leftrightarrow U)
\]

\subsection{Existence Step}\label{existence-step}

From A4:

\[
\Diamond\exists w\, Welt(w)
\]

With A2 (modal realism):

\[
\exists w\, Welt(w)
\]

Thus at least one world exists.

From the definition of \(T\):

\[
\exists w\, Welt(w) \rightarrow T
\]

Thus:

\[
\exists T
\]

With the equivalences \(T \leftrightarrow U\), \(U \leftrightarrow S\), \(S \leftrightarrow T\), it follows:

\[
\exists U \land \exists T \land \exists S
\]

Therefore:

\[
\boxed{U \land T \land S}
\]

\section{Conclusion}\label{conclusion}

We have shown:

\[
\boxed{\Box(A_1 \land A_2 \land ... \land A_{10} \rightarrow Tr)}
\]

The Trinity of Origin (U), Totality (T), and Self-Knowledge (S) follows necessarily from the axioms of a strict monistic modal ontology.

\subsection{What the proof does not show}\label{what-the-proof-does-not-show}

\begin{itemize}
\item The existence of a personal God
\item A substantial Trinity
\item A specific religious doctrine
\end{itemize}

\subsection{What the proof shows}\label{what-the-proof-shows}

\begin{itemize}
\item A certain class of monistic modal ontologies necessarily implies a Trinitarian structure.
\item The three aspects are not three substances, but three necessary perspectives on the same reality.
\end{itemize}

Thus, the Trinity reveals itself to be what philosophy has sought since Plato and Plotinus, since Augustine and Hegel: the unity that carries its difference within itself, without falling into dualism or reductionist monism.

\section{Countermodels}\label{countermodels}

The Trinity can be avoided if at least one of the axioms is abandoned:

\begin{tabular}{lll}
\textbf{Abandoned Axiom} & \textbf{Countermodel} & \textbf{Consequence} \\
A1 (Monism) & Dualism (Descartes) & S can exist without T \\
A2 (Modal Realism) & Nominalism & Possibilities need not be realized \\
A3 (No Nothing) & Eliminative Materialism & T can exist without U \\
A4 (Existence of a World) & Nihilism & There is no world \\
A5 (Consciousness) & Eliminativism & S is not necessary \\
A6 (Embedding Condition) & Contingency without Ground & U is not necessary \\
A7 (Knowledge Possibility) & Epistemic Pessimism & Complete knowledge is impossible \\
A8 (Reflexivity) & Externalism & Knowledge need not be self-referential \\
A9 (Identity of Knowledge and Object) & Epistemological Dualism & S can exist without T \\
A10 (Well-Foundedness) & Infinite Regress & There is no ultimate ground \\
\end{tabular}

\section{Comparison with Gödel's Ontological Argument}\label{comparison-with-goumldels-ontological-argument}

\subsection{Gödel's Argument (simplified)}\label{goumldels-argument-simplified}

\[
\text{Axioms about Positivity} \rightarrow \text{Necessary Existence of a Divine Being}
\]

\subsection{Comparison Table}\label{comparison-table}

\begin{tabular}{lll}
\textbf{Criterion} & \textbf{Gödel's Proof} & \textbf{This Proof} \\
Goal & Existence of a Being & Structural Necessity \\
Subject & Subject (God) & Functional Aspects (U, T, S) \\
Axioms & About "Positivity" & About the Structure of Reality \\
Formal Rigor & High (but controversial) & High (explicitly verified) \\
Metaphysical Presuppositions & Strong (concept of positivity) & Strong (monism, modal realism) \\
Philosophical Scope & Theological & Ontological-structural \\
Proximity to Classical Theology & Very high & Low (no person) \\
\end{tabular}

\subsection{Assessment}\label{assessment}

\textbf{Formal Rigor:} Both proofs are formally rigorous, but Gödel's proof suffers from the problem of "modal collapse" (if God exists necessarily, then all truths are necessary). This proof avoids this collapse because it does not postulate omnipotence or omniscience.

\textbf{Metaphysical Presuppositions:} Gödel's axioms are less intuitive (what is "positive"?). This proof uses more common-sense concepts (world, possibility, ground).

\textbf{Philosophical Scope:} Gödel's proof is more spectacular, but also more vulnerable. This proof is more modest, but also more robust.

\section{Appendix: Complete Axioms and Theorems}\label{appendix-complete-axioms-and-theorems}

\subsection{Axioms (Complete)}\label{axioms-complete}

\[
\begin{aligned}
A1 &:= \Box\neg\exists x\exists y\, FundamentalGetrennt(x,y) \\
A2 &:= \forall p(Cons(p) \rightarrow \Diamond Realisiert(p)) \land \forall p(\Diamond p \rightarrow \exists w\, Realisiert(w,p)) \\
A3 &:= \Box(N \rightarrow \neg\Diamond Welt) \\
A4 &:= \Diamond\exists w\, Welt(w) \\
A5 &:= \Diamond C \\
A6 &:= \Box(\neg Nec(p) \rightarrow \exists g\, Grund(g,p)) \\
A7 &:= \Box(C \rightarrow \Diamond V_T) \\
A8 &:= \Box(V_T(x) \rightarrow V_T(V_T(x))) \\
A9 &:= \Box(V_T(x) \rightarrow T) \\
A10 &:= \Box\neg\exists infinite\, chain(g_1, g_2, ...)
\end{aligned}
\]

\subsection{Definitions}\label{definitions}

\[
\begin{aligned}
T &:= \forall p(\Diamond p \rightarrow \exists w\, Realisiert(w,p)) \\
U &:= Grund(g, T) \land \Box \exists g \\
S &:= \forall p(K_T(p) \rightarrow K_T(K_T(p))) \\
Tr &:= U \land T \land S
\end{aligned}
\]

\subsection{Derived Theorems}\label{derived-theorems}

\[
\begin{aligned}
& \Box(T \rightarrow U) \\
& \Box(U \rightarrow S) \\
& \Box(S \rightarrow T) \\
& \Rightarrow \Box(U \leftrightarrow T) \land \Box(T \leftrightarrow S) \land \Box(S \leftrightarrow U) \\
& \Rightarrow \Box(U \land T \land S) \\
& \Rightarrow \Box Tr
\end{aligned}
\]

\section{Afterword}\label{afterword}

You have read this treatise. Perhaps the formal language remained foreign to you; perhaps you found it clarifying. Both are perfectly acceptable. For this text does not seek to proclaim a final truth, but rather to render a movement of thought comprehensible.

The Trinity of Origin, Totality, and Self-Knowledge is not an article of faith. It is the result of a consistent application of modal-logical principles to the structure of reality. Whoever shares the axioms must accept the conclusion. Whoever does not can construct countermodels.

What remains is the insight that reality, when conceived as a closed, consistent unity, necessarily exhibits three aspects: it must make itself possible, realize itself as a whole, and be able to know itself in principle. These three are not three substances, but three perspectives on one and the same reality.

Perhaps this is an answer to the question of why there is something rather than nothing. Perhaps it is merely a formally consistent description of what has always been there. The decision of what you make of it lies with you.

One final question:

Do we dare the leap into openness to the limit?

\begin{center}
\emph{This treatise was written in the spirit of strict modal logic, but in the language of philosophy -- for truth requires both: the precision of the formula and the breadth of the concept.}
\end{center}

\end{document}